\documentclass[12pt,a4paper]{article}
\usepackage[utf8]{vietnam}
\usepackage{amsmath,amssymb}
\usepackage{geometry}
\geometry{
  top=12mm,
  bottom=12mm,
  left=14mm,
  right=14mm,
  footskip=6mm,
  headheight=14pt
}

\usepackage{tikz}
\usetikzlibrary{arrows.meta,calc,shapes.geometric}
\usepackage[most]{tcolorbox}
\usepackage{fontawesome}
\usepackage{fancyhdr}
\usepackage{enumitem}

\pagestyle{fancy}
\fancyhf{}
\fancyhead[L]{\footnotesize\color{teal!80!black}\faBook\ \textbf{SỔ TAY KIẾN THỨC TOÁN 3}}
\fancyhead[R]{\footnotesize\color{teal!80!black}\textbf{Lớp Toán Thầy Lĩnh -- loptoanthaylinh.vn}}
\fancyfoot[C]{\footnotesize\color{gray}-- Sổ tay kiến thức: Bài 17. Hình tròn, tâm, bán kính, đường kính --}
\renewcommand{\headrulewidth}{0.4pt}
\renewcommand{\footrulewidth}{0pt}

\setlength{\parindent}{0pt}
\setlength{\parskip}{0pt}

\begin{document}

% --- BANNER BÀI HỌC ---
\begin{tcolorbox}[
  enhanced,
  colback=teal!8!white,
  colframe=teal!70!black,
  arc=3mm, boxrule=1.2pt,
  left=3mm, right=3mm, top=2mm, bottom=2mm,
  halign=center
]
  {\LARGE\bfseries\color{teal!90!black}\faBookmark\ BÀI 17: HÌNH TRÒN, TÂM, BÁN KÍNH, ĐƯỜNG KÍNH}\\[1mm]
  {\itshape\footnotesize\color{black!75} Giải phẫu cấu trúc hình tròn, quan hệ đường kính gấp đôi bán kính, kĩ thuật vẽ compa trên 1 trang A4}
\end{tcolorbox}

\vspace{-2mm}

\noindent
\begin{minipage}[t]{0.49\textwidth}\vspace{0pt}
\begin{tcolorbox}[
  enhanced,
  colback=blue!4!white,
  colframe=blue!75!black,
  arc=2mm, boxrule=0.8pt,
  left=2mm, right=2mm, top=1.8mm, bottom=1.8mm,
  title={\footnotesize\bfseries\color{white}\faDotCircleO\ 1. Tâm, Bán kính \& Đường kính}
]
{\footnotesize
\begin{itemize}\setlength{\itemsep}{1pt}\setlength{\leftskip}{-2.5mm}
  \item \textbf{Tâm $O$:} Điểm chính giữa hình tròn.
  \item \textbf{Bán kính ($R$):} Đoạn thẳng nối tâm $O$ với một điểm bất kì trên đường tròn (VD: $OA, OB, OM$).
  \begin{itemize}\setlength{\itemsep}{0.5pt}\setlength{\leftskip}{-2.5mm}
    \item Tất cả các bán kính đều \textbf{bằng nhau}.
  \end{itemize}
  \item \textbf{Đường kính ($d$):} Đoạn thẳng nối hai điểm trên đường tròn và \textbf{đi qua tâm $O$} (VD: đoạn thẳng $AB$).
  \begin{itemize}\setlength{\itemsep}{0.5pt}\setlength{\leftskip}{-2.5mm}
    \item Tâm $O$ chính là \textbf{trung điểm} của đường kính $AB$.
  \end{itemize}
\end{itemize}
}
\end{tcolorbox}
\end{minipage}
\hfill
\begin{minipage}[t]{0.49\textwidth}\vspace{0pt}
\begin{tcolorbox}[
  enhanced,
  colback=teal!4!white,
  colframe=teal!75!black,
  arc=2mm, boxrule=0.8pt,
  left=2mm, right=2mm, top=1.8mm, bottom=1.8mm,
  title={\footnotesize\bfseries\color{white}\faBalanceScale\ 2. Mối quan hệ \& Cách vẽ compa}
]
{\footnotesize
\begin{itemize}\setlength{\itemsep}{1pt}\setlength{\leftskip}{-2.5mm}
  \item \textbf{Quy tắc vàng:} Đường kính dài \textbf{gấp 2 lần} bán kính (Bán kính bằng \textbf{một nửa} đường kính):
  $$\mathbf{d = 2 \times R} \iff \mathbf{R = d : 2}$$
  \item \textbf{Cách vẽ hình tròn bằng com-pa:}
  \begin{enumerate}\setlength{\itemsep}{0.5pt}\setlength{\leftskip}{-2.5mm}
    \item Mở rộng com-pa bằng đúng độ dài bán kính $R$ trên thước kẻ.
    \item Đặt đầu nhọn của com-pa cố định vào vị trí tâm $O$.
    \item Quay đầu bút chì $1$ vòng quanh tâm.
  \end{enumerate}
\end{itemize}
}
\end{tcolorbox}
\end{minipage}

\vspace{-2.5mm}
% --- KHỐI 3: SƠ ĐỒ TƯ DUY TIKZ ---
\noindent
\begin{tcolorbox}[
  enhanced,
  colback=orange!4!white,
  colframe=orange!80!black,
  arc=2mm, boxrule=0.8pt,
  left=2mm, right=2mm, top=1.5mm, bottom=1.5mm,
  title={\footnotesize\bfseries\color{white}\faSitemap\ SƠ ĐỒ TƯ DUY: CẤU TRÚC HÌNH HỌC HÌNH TRÒN TÂM $O$}
]
\centering
\begin{tikzpicture}[scale=0.8]
  \draw[thick, blue!80!black, fill=blue!5] (0, 0) circle (1.5cm);
  \coordinate (O) at (0, 0);
  \coordinate (A) at (-1.5, 0);
  \coordinate (B) at (1.5, 0);
  \coordinate (C) at (0, 1.5);

  \draw[thick, red!80!black] (A) -- (B) node[midway, below, font=\scriptsize\bfseries] {Đường kính $AB = 2R$};
  \draw[thick, teal!80!black] (O) -- (C) node[midway, right, font=\scriptsize\bfseries] {Bán kính $R$};

  \fill[black] (O) circle (1.5pt) node[below right, font=\scriptsize\bfseries] {$O$};
  \fill[blue!80!black] (A) circle (1.5pt) node[left, font=\scriptsize\bfseries] {$A$};
  \fill[blue!80!black] (B) circle (1.5pt) node[right, font=\scriptsize\bfseries] {$B$};
  \fill[blue!80!black] (C) circle (1.5pt) node[above, font=\scriptsize\bfseries] {$C$};
\end{tikzpicture}
\end{tcolorbox}

\vspace{-2.5mm}
% --- KHỐI 4: VÍ DỤ & MẸO TRÁNH BẪY ---
\noindent
\begin{minipage}[t]{0.49\textwidth}\vspace{0pt}
\begin{tcolorbox}[
  enhanced,
  colback=green!4!white,
  colframe=green!60!black,
  arc=2mm, boxrule=0.8pt,
  left=2mm, right=2mm, top=1.8mm, bottom=1.8mm,
  title={\footnotesize\bfseries\color{white}\faCheckCircle\ Ví dụ thực hành chuẩn mực}
]
{\footnotesize
\textbf{1. Cho bán kính tìm đường kính:}\\
Hình tròn có bán kính $R = 4\text{ cm}$.\\
Đường kính là: $4 \times 2 = \mathbf{8\text{ cm}}$.\\
\textbf{2. Cho đường kính tìm bán kính:}\\
Hình tròn có đường kính $d = 12\text{ cm}$.\\
Bán kính là: $12 : 2 = \mathbf{6\text{ cm}}$.\\
\textbf{3. Mở compa vẽ hình:}\\
Muốn vẽ hình tròn đường kính $10\text{ cm}$, ta phải mở compa rộng bằng bán kính là $10 : 2 = \mathbf{5\text{ cm}}$.
}
\end{tcolorbox}
\end{minipage}
\hfill
\begin{minipage}[t]{0.49\textwidth}\vspace{0pt}
\begin{tcolorbox}[
  enhanced,
  colback=red!4!white,
  colframe=red!70!black,
  arc=2mm, boxrule=0.8pt,
  left=2mm, right=2mm, top=1.8mm, bottom=1.8mm,
  title={\footnotesize\bfseries\color{white}\faExclamationTriangle\ Mẹo tránh bẫy trắc nghiệm}
]
{\footnotesize
\begin{itemize}\setlength{\itemsep}{1.5pt}\setlength{\leftskip}{-2.5mm}
  \item \textbf{Bẫy đường kính không đi qua tâm:}\\
  Đoạn thẳng nối 2 điểm trên hình tròn nhưng KHÔNG đi qua tâm thì KHÔNG PHẢI là đường kính!
  \item \textbf{Bẫy mở compa bằng đường kính:}\\
  Đề bài yêu cầu vẽ hình tròn đường kính $8\text{ cm}$, nhiều bạn mở compa $8\text{ cm}$ (SAI!). Độ mở compa luôn luôn bằng \textbf{bán kính} ($8 : 2 = 4\text{ cm}$).
\end{itemize}
}
\end{tcolorbox}
\end{minipage}


\end{document}
